\documentclass[10pt,a4paper]{article}

\usepackage[a4paper, left=20mm, right=20mm, top=15mm, bottom=15mm]{geometry}
\usepackage{fontspec}
\setmainfont{TeX Gyre Heros}
\usepackage{amsmath,amssymb,mathtools}
\usepackage{array,booktabs}
\usepackage{enumitem}

\setlength{\parskip}{3pt}
\setlength{\parindent}{0pt}
\setlength{\abovedisplayskip}{4pt}
\setlength{\belowdisplayskip}{4pt}
\setlength{\abovedisplayshortskip}{2pt}
\setlength{\belowdisplayshortskip}{2pt}

\pagestyle{plain}

\newcommand{\sectitle}[1]{%
  \vspace{8pt}\noindent\rule{\linewidth}{0.8pt}\par\vspace{2pt}%
  {\large\textbf{#1}}\par\vspace{2pt}\rule{\linewidth}{0.4pt}\vspace{4pt}}

\newcommand{\Sol}{\par\vspace{4pt}\noindent\textbf{Solution:}\enspace}

\usepackage{eso-pic}
\usepackage{tikz}
\usetikzlibrary{calc}
\usepackage[useregional]{datetime2}

% Running margin label
\newcommand{\mymarginlabel}{\textcopyright\ 2026 David Ewing dew59@uclive.ac.nz | \DTMnow\ | \jobname.tex}
\AddToShipoutPictureBG{%
  \begin{tikzpicture}[remember picture,overlay]
    \node[rotate=90, anchor=north]
      at ($(current page.north west)+(1.4cm,-13cm)$)
      {\scriptsize\ttfamily \mymarginlabel};
  \end{tikzpicture}%
}

\begin{document}

\begin{center}
  {\large\textbf{ENEL445 Mock Test --- Questions \& Solutions}}\\[2pt]
  {\small Based on 2024 mid-year exam and 2025 mid-year exam. Questions and mark
  allocations taken directly from past papers.}
\end{center}

\noindent\rule{\linewidth}{0.4pt}

% ==============================================================
\sectitle{Q1. Gradient-Free Optimisation (Part I)}
% ==============================================================

\noindent\textbf{(a)} [3 marks]\quad Give three situations where gradient-free methods are more suited than gradient-based methods for solving an optimisation problem.

\Sol
\begin{enumerate}[noitemsep, topsep=2pt]
  \item \textbf{Non-differentiable or discontinuous} objective --- gradient does
        not exist or is unreliable (e.g.\ simulation outputs, physical
        experiments, $|x|$ at $x=0$).
  \item \textbf{Black-box function} --- no analytical expression is available;
        only function evaluations are possible (e.g.\ legacy simulation code,
        hardware-in-the-loop testing).
  \item \textbf{Noisy function evaluations} --- gradient estimates are corrupted
        by noise, making gradient-based methods unreliable or divergent.
\end{enumerate}

\vspace{4pt}
\noindent\textbf{(b)} [2 marks]\quad Describe (or sketch) the Nelder-Mead
\textit{expansion} operation.

\Sol The worst vertex $x_w$ is reflected through the centroid $\bar{x}$ of the
remaining vertices:
\[
  x_r = \bar{x} + (\bar{x} - x_w) \quad\text{(reflection)}.
\]
Expansion then extends further:
\[
  x_e = \bar{x} + 2(\bar{x} - x_w).
\]
If $f(x_e) < f(x_r)$, accept $x_e$; otherwise accept $x_r$. On a contour plot
$x_e$ lies on the opposite side of $\bar{x}$ from $x_w$, twice as far as the
reflection point.

\vspace{4pt}
\noindent\textbf{(c)} [3 marks]\quad Name three convergence metrics used with
Nelder-Mead.

\Sol
\begin{enumerate}[noitemsep, topsep=2pt]
  \item \textbf{Simplex size}: $\max_i \|x_i - x_1\|_2 < \varepsilon$
  \item \textbf{Function value spread}: $\max_i |f(x_i) - f(x_1)| < \varepsilon$
  \item \textbf{Iteration / evaluation limit}: $k > k_{\max}$ or number of
        function evaluations exceeds $N_{\max}$
\end{enumerate}

\vspace{4pt}
\noindent\textbf{(d)} [2 marks]\quad Sketch a minimal positive spanning set in
$\mathbb{R}^2$ \textbf{and} state the process of opportunistic polling in the
generalised pattern search algorithm.

\Sol A minimal positive spanning set in $\mathbb{R}^2$ requires $n+1=3$ vectors.
Example: $\{e_1,\, e_2,\, -e_1-e_2\}$ where $e_1=(1,0)^\top$,
$e_2=(0,1)^\top$. These span all of $\mathbb{R}^2$ in the positive sense.

\textbf{Opportunistic polling:} poll directions are evaluated one at a time in a
preferred order. As soon as one direction yields $f < f_{\text{current}}$, the
algorithm immediately accepts that point and starts a new iteration --- it does
\emph{not} evaluate all remaining directions in the pattern.

\vspace{4pt}
\noindent\textbf{(e)} [3 marks]\quad State the two significant shortcomings of
Shubert's algorithm and briefly explain how DIRECT addresses them.

\Sol
\textbf{Shubert shortcomings:}
\begin{enumerate}[noitemsep, topsep=2pt]
  \item Requires the global Lipschitz constant $K$ to be known in advance ---
        which is generally unknown in practice.
  \item Only applicable to \emph{one-dimensional} problems.
\end{enumerate}

\textbf{DIRECT addresses both:}
\begin{enumerate}[noitemsep, topsep=2pt]
  \item Estimates local Lipschitz constants implicitly per hyper-rectangle by
        using the $f$-$d$ plot --- no global $K$ is needed.
  \item Generalises to $n$ dimensions by partitioning $n$-dimensional
        hyper-rectangles.
\end{enumerate}

\vspace{4pt}
\noindent\textbf{(f)} [3 marks]\quad On the $f$-$d$ plot, identify which
rectangles DIRECT subdivides next when $\varepsilon|f_{\min}|=3$.

\Sol DIRECT selects all rectangles whose sample point $(d_i, f(c_i))$ lies on
the \emph{lower-right convex hull} of the $f$-$d$ scatter and satisfies:
\[
  f(c_i) - K d_i \leq f_{\min} - \varepsilon|f_{\min}| \quad \text{for some } K > 0.
\]
With $\varepsilon|f_{\min}|=3$, points are eligible only if their value is within
3 of $f_{\min}$ after accounting for the hull slope. On the plot: circle all
points on the lower-right convex hull (points not dominated below-right by any
other point).

\vspace{4pt}
\noindent\textbf{(g)} [4 marks]\quad Briefly explain the three main steps in a
genetic algorithm and discuss the advantages and disadvantages of a binary-coded
gene versus a real-coded gene.

\Sol
\textbf{Three steps:}
\begin{enumerate}[noitemsep, topsep=2pt]
  \item \textbf{Selection}: choose parents with probability proportional to
        fitness (e.g.\ roulette-wheel or tournament selection).
  \item \textbf{Crossover}: combine genetic material of two parents at a random
        cut point to produce offspring (e.g.\ single-point crossover swaps tails).
  \item \textbf{Mutation}: randomly flip/alter genes with a small probability to
        maintain diversity and prevent premature convergence to a local optimum.
\end{enumerate}

\textbf{Binary vs real-coded:}
\begin{itemize}[noitemsep, topsep=2pt]
  \item \textbf{Binary}: simple, well-studied operators; but suffers from the
        \emph{Hamming cliff} --- numerically adjacent integers can differ in many
        bits, causing large jumps and disrupting convergence on smooth problems.
  \item \textbf{Real-coded}: natural for continuous variables; no Hamming cliff;
        precision is not limited by bit length; better suited to smooth continuous
        optimisation.
\end{itemize}

% ==============================================================
\sectitle{Q1 (2025 variant). Additional Gradient-Free Questions}
% ==============================================================

\noindent\textbf{(h)} [3 marks]\quad The particle swarm update is:
\[
  x^{(i)}_{k+1} = x^{(i)}_k + v^{(i)}_{k+1}\Delta t, \qquad
  v^{(i)}_{k+1} = \alpha v^{(i)}_k
    + \beta\frac{x^{(i)}_{\mathrm{best}}-x^{(i)}_k}{\Delta t}
    + \gamma\frac{x_{\mathrm{best}}-x^{(i)}_k}{\Delta t}.
\]
Define all symbols.

\Sol
\begin{itemize}[noitemsep, topsep=2pt]
  \item $x^{(i)}_k$: position of particle $i$ at iteration $k$
  \item $v^{(i)}_k$: velocity of particle $i$ at iteration $k$
  \item $x^{(i)}_{\mathrm{best}}$: best position ever visited by particle $i$
        (personal best)
  \item $x_{\mathrm{best}}$: best position ever found by \emph{any} particle
        in the swarm (global best)
  \item $\alpha$: inertia weight (momentum of previous velocity)
  \item $\beta$: cognitive coefficient (attraction to personal best)
  \item $\gamma$: social coefficient (attraction to global best)
  \item $\Delta t$: time step; $k$ = iteration index; $i$ = particle index
\end{itemize}

% ==============================================================
\sectitle{Q2. Gradient-Free Optimisation (Part II)}
% ==============================================================

\noindent\textbf{(a)} [5 marks]\quad Convert to standard LP form
(minimise, equality constraints, $x\geq0$):
\[
  \max\ 2x_1 - x_2 + 8x_3 \quad\text{s.t.}\quad
  7x_1+5x_2-x_3\leq1,\quad x_1+6x_2\geq4,\quad
  x_1,x_2\geq0,\ x_3\ \text{free}.
\]

\Sol Four steps:
\begin{enumerate}[noitemsep, topsep=2pt]
  \item \textbf{Objective}: negate $\;\Rightarrow\;$ $\min\ {-2x_1+x_2-8x_3}$
  \item \textbf{Free variable}: $x_3 = x_3^+ - x_3^-$,\; $x_3^+,x_3^-\geq0$
  \item \textbf{$\leq$ constraint}: add slack $s_1\geq0$:\;
        $7x_1+5x_2-x_3^++x_3^-+s_1=1$
  \item \textbf{$\geq$ constraint}: subtract surplus $s_2\geq0$:\;
        $x_1+6x_2-s_2=4$
\end{enumerate}

\textbf{Standard form:}
\[
  \min\ -2x_1+x_2-8x_3^++8x_3^-
  \quad\text{s.t.}\quad
  \underbrace{\begin{bmatrix}7&5&-1&1&1&0\\1&6&0&0&0&-1\end{bmatrix}}_{A}
  \begin{bmatrix}x_1\\x_2\\x_3^+\\x_3^-\\s_1\\s_2\end{bmatrix}
  =\begin{bmatrix}1\\4\end{bmatrix},\quad
  \text{all variables}\geq0.
\]

\vspace{4pt}
\noindent\textbf{(b)} [4 marks]\quad Solve the following optimal assignment
problem using the Hungarian algorithm (find the assignment with minimum cost):

\begin{center}
\begin{tabular}{lccc}
\toprule
 & Task 1 & Task 2 & Task 3 \\
\midrule
Worker A & 100 & 150 & 125 \\
Worker B & 150 & 175 & 135 \\
Worker C & 120 & 250 & 140 \\
\bottomrule
\end{tabular}
\end{center}

\Sol
\textbf{Step 1 --- row reduction} (subtract row minimum: A=100, B=135, C=120):
\[
  \begin{pmatrix}0&50&25\\15&40&0\\0&130&20\end{pmatrix}
\]

\textbf{Step 2 --- column reduction} (subtract column minimum: col1=0, col2=40, col3=0):
\[
  \begin{pmatrix}0&10&25\\15&0&0\\0&90&20\end{pmatrix}
\]

\textbf{Step 3 --- cover zeros}: zeros at $(A,T_1),(B,T_2),(B,T_3),(C,T_1)$.
Minimum cover: col~$T_1$ + row~$B$ = 2 lines $< n=3$. Not optimal.

\textbf{Step 4 --- update}: min uncovered element =
$\min\{10,25,90,20\}=10$. Subtract 10 from uncovered cells; add 10 to
doubly-covered cell $(B,T_1)$:
\[
  \begin{pmatrix}0&0&15\\25&0&0\\0&80&10\end{pmatrix}
\]

\textbf{Step 5}: zeros at $(A,T_1),(A,T_2),(B,T_2),(B,T_3),(C,T_1)$.
Cover with row~$A$ + row~$B$ + col~$T_1$ = 3 lines $= n$. Optimal.

\textbf{Step 6 --- assignment} (independent zeros):
$B\!\to\! T_3$\; (unique zero in $T_3$),\; $C\!\to\! T_1$,\; $A\!\to\! T_2$.

\textbf{Total cost:} $150 + 135 + 120 = \mathbf{405}$.

\vspace{4pt}
\noindent\textbf{(c)} [2 marks]\quad State the purpose of the relaxation step
in branch-and-bound \textbf{and} give an example.

\Sol \textbf{Purpose:} solve the LP relaxation (drop integrality constraints)
to obtain a bound on the optimal integer solution. If the relaxation bound is
worse than the current best integer solution, the entire branch is pruned.

\textbf{Example:} in a binary programme with $x_i\in\{0,1\}$, relax to
$0\leq x_i\leq1$ (continuous LP). Solve; if $x_i^*\notin\{0,1\}$, branch
by adding $x_i\leq0$ and $x_i\geq1$ as two sub-problems.

\vspace{4pt}
\noindent\textbf{(d)} [3 marks]\quad Define the Pareto front and explain how it
provides insights into an optimisation problem.

\Sol A solution $x^*$ is \textbf{Pareto-optimal} if no other feasible solution
simultaneously improves one objective without worsening at least one other. The
\textbf{Pareto front} is the set of all such non-dominated solutions.

\textbf{Insights:} the front reveals the trade-off structure between competing
objectives. A decision-maker can inspect the shape of the front to quantify
trade-offs (e.g.\ how much cost increases per unit reduction in emissions) and
select a preferred solution. A single-objective optimum lies at one extreme of
the front.

\vspace{4pt}
\noindent\textbf{(e)} [6 marks]\quad Show that maximising the log-likelihood is
equivalent to minimising the sum of squared errors when measurement errors are
i.i.d.\ $\mathcal{N}(0,\sigma^2)$.

\Sol Let $z_n = g_n - m(c,y_n)$ be the error at measurement $n$, distributed as
$z_n\sim\mathcal{N}(0,\sigma^2)$. The likelihood is:
\[
  L(c) = \prod_{n=1}^{N} \frac{1}{\sqrt{2\pi\sigma^2}}
         \exp\!\left(-\frac{z_n^2}{2\sigma^2}\right).
\]
Taking the log:
\[
  \ln L(c) = -\frac{N}{2}\ln(2\pi\sigma^2)
             - \frac{1}{2\sigma^2}\sum_{n=1}^{N} z_n^2.
\]
The first term is a constant (independent of $c$); $\frac{1}{2\sigma^2}>0$ is a
positive constant. Therefore maximising $\ln L$ over $c$ is equivalent to
minimising $\sum_{n=1}^N z_n^2$. $\blacksquare$

\vspace{4pt}
\noindent\textbf{(2025 variant: simulated annealing)} [3 marks]\quad Briefly
explain simulated annealing, incorporating the terms \textit{acceptance
probability} and \textit{cooling schedule}.

\Sol Simulated annealing is a stochastic local search. At each iteration a
candidate $x'$ is generated near the current point $x$:
\begin{itemize}[noitemsep, topsep=2pt]
  \item If $f(x') < f(x)$: accept unconditionally.
  \item If $f(x') \geq f(x)$: accept with \textbf{acceptance probability}
        $P = \exp\!\bigl(-(f(x')-f(x))/T\bigr)$, where $T$ is the current
        temperature.
\end{itemize}
The \textbf{cooling schedule} reduces $T$ over time (e.g.\ $T_k = T_0\cdot r^k$
for $r<1$). High $T$ allows uphill moves (exploration); as $T\to0$ the algorithm
converges to pure descent (exploitation).

% ==============================================================
\sectitle{Q3. Unconstrained Optimisation}
% ==============================================================

\noindent\textbf{(a)(i)} [2 marks]\quad Represent $i=1,2,3,4$ in base $b=3$.

\Sol Using $i = a_0 + a_1\cdot3 + a_2\cdot9+\cdots$:
\[
  1=(1)_3, \quad 2=(2)_3, \quad 3=(10)_3, \quad 4=(11)_3.
\]

\vspace{4pt}
\noindent\textbf{(a)(ii)} [2 marks]\quad Compute $\phi_i(3)$ for $i=1,2,3,4$
using $\phi_i(b)=\frac{a_0}{b}+\frac{a_1}{b^2}+\cdots$

\Sol
\[
  \phi_1(3)=\tfrac{1}{3},\quad
  \phi_2(3)=\tfrac{2}{3},\quad
  \phi_3(3)=\tfrac{0}{3}+\tfrac{1}{9}=\tfrac{1}{9},\quad
  \phi_4(3)=\tfrac{1}{3}+\tfrac{1}{9}=\tfrac{4}{9}.
\]

\vspace{4pt}
\noindent\textbf{(b)(i)} [3 marks]\quad Using the second-order Taylor expansion
of $f(x)$ around $x_{k-1}$, find the Newton update $x_k$.

\Sol Minimise the approximation $q(x)$ by differentiating and setting to zero:
\[
  \frac{\partial q}{\partial x}
  = \nabla f(x_{k-1}) + H(x_{k-1})(x-x_{k-1}) = 0
  \;\Rightarrow\;
  \boxed{x_k = x_{k-1} - H(x_{k-1})^{-1}\nabla f(x_{k-1})}.
\]

\vspace{4pt}
\noindent\textbf{(b)(ii)} [4 marks]\quad Show that $(x_k-x_{k-1})$ lies in the
direction of decreasing $f$ when $H(x_{k-1})$ is positive definite.

\Sol Let $d = x_k - x_{k-1} = -H^{-1}\nabla f$. A direction is a descent
direction if $\nabla f^\top d < 0$:
\[
  \nabla f^\top d
  = \nabla f^\top(-H^{-1}\nabla f)
  = -\nabla f^\top H^{-1}\nabla f.
\]
Since $H$ is positive definite (PD), $H^{-1}$ is also PD, so
$\nabla f^\top H^{-1}\nabla f > 0$ for $\nabla f\neq0$. Therefore
$\nabla f^\top d < 0$, confirming $d$ is a descent direction. $\blacksquare$

\vspace{4pt}
\noindent\textbf{(b)(iii)} [4 marks]\quad Using the quasi-Newton approximation
$f(x)\approx f(x_k)+\nabla f(x_k)^\top(x-x_k)+\tfrac{1}{2}(x-x_k)^\top
\tilde{H}(x_k)(x-x_k)$, write out the secant condition.

\Sol Differentiate the approximation with respect to $x$:
\[
  \nabla f(x) \approx \nabla f(x_k) + \tilde{H}(x_k)(x - x_k).
\]
Set $x = x_{k-1}$:
\[
  \nabla f(x_{k-1}) = \nabla f(x_k) + \tilde{H}(x_k)(x_{k-1}-x_k).
\]
Define $s_{k-1}=x_k-x_{k-1}$ and $y_{k-1}=\nabla f(x_k)-\nabla f(x_{k-1})$.
Rearranging:
\[
  \boxed{\tilde{H}(x_k)\,s_{k-1} = y_{k-1}.}
\]
This is the secant condition: the approximate Hessian must reproduce the
observed change in gradient over the last step.

\vspace{4pt}
\noindent\textbf{(c)} [5 marks]\quad Find $x^*$ minimising the ridge regression
objective $f(x)=\tfrac{1}{2}(y-Ax)^\top(y-Ax)+\tfrac{\lambda}{2}x^\top x$,
$\lambda>0$.

\Sol Differentiate and set to zero:
\[
  \nabla_x f = -A^\top(y-Ax) + \lambda x = A^\top Ax - A^\top y + \lambda x = 0.
\]
\[
  (A^\top A + \lambda I)\,x = A^\top y
  \;\Rightarrow\;
  \boxed{x^* = (A^\top A + \lambda I)^{-1}A^\top y.}
\]
Since $\lambda>0$, $(A^\top A+\lambda I)$ is positive definite, hence always
invertible, giving a unique global minimiser.

\vspace{6pt}
\noindent\rule{\linewidth}{0.4pt}
\vspace{2pt}

\noindent\textbf{2025 Q3b variant} --- the objective is:
$L(x) = f(x_0)+\nabla f(x_0)^\top(x-x_0)+\tfrac{1}{2\mu}(x-x_0)^\top(x-x_0)$

\textbf{(i)} Gradient and Hessian:
\[
  \nabla_x L = \nabla f(x_0) + \tfrac{1}{\mu}(x-x_0),
  \qquad
  H_L = \tfrac{1}{\mu}I.
\]

\textbf{(ii)} Set $\nabla_x L=0$:\;
$x^* = x_0 - \mu\,\nabla f(x_0)$ (gradient descent step with step size $\mu$).
Local minimum since $H_L = \tfrac{1}{\mu}I \succ 0$ for $\mu>0$.

\textbf{(iii)} If $\mu<0$: $H_L = \tfrac{1}{\mu}I \prec 0$ (negative definite)
--- $L$ is a concave bowl opening downward with no lower bound.
$\|x\|_2\to\infty$; no bounded solution exists.

% ==============================================================
\sectitle{Q4. Power Systems Optimisation}
% ==============================================================

\noindent\textbf{(a)} [4 marks]\quad Estimate the time to solve AC OPF given:
DC OPF solves in 5 s; 20 generators, 1000 buses, 1000 lines;
$O((n+n_g+n_h)^3)$ scaling; nonlinearity gives a $5\times$ slowdown.

\Sol \textbf{DC OPF:} $\sim n=1000$ bus angles, $n_h\sim1000$ (power balance),
$n_g\sim1000$ (line limits). Complexity $\propto(n+n_g+n_h)^3=2000^3$, time = 5 s.

\textbf{AC OPF:} add voltage magnitudes and reactive power --- roughly doubles
variables and constraints $\Rightarrow$ $n+n_g+n_h\approx4000$:
\[
  t_{\text{AC}} = 5\,\text{s} \times \frac{4000^3}{2000^3} \times 5 = 5\times8\times5 = 200\,\text{s}.
\]

\vspace{4pt}
\noindent\textbf{(b)} [6 marks]\quad Suggest an algorithm for AC OPF and
explain its main concepts.

\Sol \textbf{Algorithm: Interior Point Method (IPM)}.

\textbf{Justification:} AC OPF is a large-scale non-convex NLP with both
equality and inequality constraints. IPMs are efficient for large sparse
structured problems (e.g.\ IPOPT used in industry).

\textbf{Key concepts:} convert inequality constraints $g(x)\leq0$ to equality
using slack variables and add a logarithmic barrier:
\[
  \min_{x,s}\ f(x) - \mu\sum_i\ln s_i
  \quad\text{s.t.}\quad h(x)=0,\; g(x)+s=0,\; s>0.
\]
As $\mu\to0$ the barrier solution converges to the constrained optimum. Each
iteration solves a linear KKT system (Newton step on the first-order conditions):
\[
  \begin{bmatrix}H_{\mathcal{L}} & J_h^\top & J_g^\top\\
  J_h & 0 & 0\\ J_g & 0 & -S^{-1}Z\end{bmatrix}
  \begin{bmatrix}\Delta x\\\Delta\nu\\\Delta\lambda\end{bmatrix}
  = -r,
\]
exploiting sparsity of the power network for efficiency.

\vspace{4pt}
\noindent\textbf{(c)} [5 marks]\quad What is the difference between SC-OPF and
regular AC OPF? Why is SC-OPF hard to solve in practice?

\Sol \textbf{Regular AC OPF:} minimise generation cost subject to power flow
equations, generator limits, and thermal line limits for the base (intact)
operating case only.

\textbf{SC-OPF (N-1 security constrained):} additionally requires the system
to remain feasible after any single component outage (any one line or
generator failure). Adds $N_{\text{components}}$ full sets of contingency
constraints.

\textbf{Why hard:} for a 1000-bus, 1000-line system this means $\sim10^3$
additional complete AC OPF problems. Direct formulation has $O(N_{\text{lines}})$
times more constraints, making it computationally intractable. Non-convexity
also means each sub-problem may not solve to global optimality.

\vspace{4pt}
\noindent\textbf{(d)} [2 marks]\quad How are security constraints included in
industry practice without incurring huge computational costs?

\Sol \textbf{Contingency screening:} quickly rank all $N$ contingency cases
using a fast DC approximation. Select only the small subset of
\textit{binding} (worst-case) contingencies, then resolve AC OPF with
those constraints added. Repeat iteratively until no new binding
contingencies are found. This avoids solving the full SC-OPF directly.

\vspace{4pt}
\noindent\textbf{(e)} [8 marks]\quad Formulate a non-linear least squares
problem to estimate power system states from measurements
$P_{12},Q_{12},P_{13},Q_{13},P_{32},Q_{32},V_1,V_2$, where:
\[
  P_{ij} = V_iV_jB_{ij}\sin(\theta_i-\theta_j),\qquad
  Q_{ij} = V_i(V_i-V_j\cos(\theta_i-\theta_j))B_{ij}.
\]

\Sol \textbf{States:} set $\theta_1=0$ as reference. Unknown vector:
$x = [\theta_2,\,\theta_3,\,V_1,\,V_2,\,V_3]^\top$ (5 unknowns).

\textbf{Measurement model} $h(x)$:
\[
  h(x) = \begin{bmatrix}
    V_1V_2B_{12}\sin(\theta_1-\theta_2)\\
    V_1(V_1-V_2\cos(\theta_1-\theta_2))B_{12}\\
    V_1V_3B_{13}\sin(\theta_1-\theta_3)\\
    V_1(V_1-V_3\cos(\theta_1-\theta_3))B_{13}\\
    V_3V_2B_{32}\sin(\theta_3-\theta_2)\\
    V_3(V_3-V_2\cos(\theta_3-\theta_2))B_{32}\\
    V_1\\ V_2
  \end{bmatrix},\quad
  z = \begin{bmatrix}P_{12}\\Q_{12}\\P_{13}\\Q_{13}\\P_{32}\\Q_{32}\\V_1\\V_2\end{bmatrix}.
\]

\textbf{Optimisation problem} (all meters equal variance $\sigma^2$):
\[
  \min_{x}\; \|z - h(x)\|_2^2
  = \min_{x}\;\sum_{m=1}^{8}(z_m - h_m(x))^2.
\]
Solved numerically using Gauss-Newton or Levenberg-Marquardt iteration.

\end{document}
